% !TEX root = web.tex

% Web variant of the blueprint, used by `leanblueprint` to produce the
% online HTML blueprint via plastex. Loads the same content but with the
% `blueprint` package which provides clickable Lean/dependency cross-links.

\documentclass{report}

\usepackage[utf8]{inputenc}
\usepackage{amsmath, amssymb, amsthm}
\usepackage{hyperref}
\usepackage[capitalize]{cleveref}

% Loads the real macros for \lean, \leanok, \uses, \notready, \proves
\usepackage{blueprint}

\newtheorem{theorem}{Theorem}[chapter]
\newtheorem{proposition}[theorem]{Proposition}
\newtheorem{lemma}[theorem]{Lemma}
\newtheorem{corollary}[theorem]{Corollary}
\newtheorem{remark}[theorem]{Remark}
\theoremstyle{definition}
\newtheorem{definition}[theorem]{Definition}
\newtheorem{example}[theorem]{Example}

% plasTeX renders thebibliography as a plain list with no heading, which leaves
% the references glued to the end of the last section. A chapter of their own
% gives them a heading, their own page (split-level 2) and an entry in the table
% of contents. It carries no \label: plasTeX binds a label to the last
% *numbered* node (Context.label via refstepcounter), so a label here would
% attach to the final theorem instead. The page is therefore named
% sect0001.html rather than after a label.
% The PDF build defines this differently -- see print.tex.
\newcommand{\bpbibheading}{\chapter*{References}}

\title{Blueprint:\\ A Borsuk--Ulam Theorem for\\
       Adaptive Continuous Measurements}
\author{Mario Ullrich}

\begin{document}
\maketitle
% Links rendered below the title: the GitHub Pages deployment and the repository.
\home{https://mario-ullrich.github.io/Lean-ContMeasurements/}
\github{https://github.com/mario-ullrich/Lean-ContMeasurements}

% !TEX root = print.tex

\chapter{Setting and the model}
\label{ch:model}

Throughout, $E$ is a real normed space, finite-dimensional where a dimension is
named, and $d = \dim E$. We write $S_R = \{x \in E : \|x\| = R\}$ for the sphere of
radius $R$ and $B_R$ for the closed ball.

A point $x \in E$ is to be recovered from $n$ real numbers, each of which may be
chosen after the previous ones have been seen. Every single measurement is continuous
in $x$; the rule choosing it from the observed values is arbitrary. The resulting
information map is therefore in general discontinuous, which is what lets $n$ be
exponentially smaller than $d$.

Each statement below names the Lean declaration that formalises it. Every statement
and every proof is formalised: the project compiles without \texttt{sorry}, so each
result is verified by the Lean kernel.

\begin{definition}[Adaptive scheme]
\label{def:adaptive-scheme}
\lean{ContMeasurements.AdaptiveScheme}
\leanok
An \emph{adaptive scheme} on a topological space $X$ is a family
\[
  \lambda_k^{(y)} : X \longrightarrow \mathbb{R},
  \qquad k \in \mathbb{N}_0, \quad y \in \mathbb{R}^k,
\]
of continuous functionals, one for every number $k$ of values already observed and
every such vector of values.
\end{definition}

The functional asked first is $\lambda_1 = \lambda_0^{()}$, and after observing the
value $y$ the scheme continues with the family $\lambda_{k+1}^{(y, \cdot)}$; these are
\texttt{first} and \texttt{tail} in the formalisation.

\begin{definition}[Information map]
\label{def:info}
\lean{ContMeasurements.AdaptiveScheme.info}
\leanok
\uses{def:adaptive-scheme}
For an adaptive scheme $S$ on $X$ and $n \in \mathbb{N}_0$, the \emph{information map}
$N_n : X \to \mathbb{R}^n$ is defined by recursion on $n$: $N_0(x)$ is the empty
tuple, and
\[
  N_{n+1}(x) = \bigl( \lambda_1(x), \, N_n'(x) \bigr),
\]
where $N_n'$ is the information map of the scheme that continues after the value
$\lambda_1(x)$.
\end{definition}

An \emph{algorithm} is $\Phi \circ N_n$ with an arbitrary map
$\Phi : \mathbb{R}^n \to X$. Errors are stated pointwise rather than through a
supremum, which avoids side conditions and says the same thing.

\begin{definition}[Distinguishing antipodes]
\label{def:distinguishes-antipodes}
\lean{ContMeasurements.AdaptiveScheme.DistinguishesAntipodes}
\leanok
\uses{def:info}
A scheme $S$ \emph{distinguishes antipodes} on a set $K \subseteq X$ with $n$
measurements if $N_n(x) \neq N_n(-x)$ for every $x \in K$.
\end{definition}

The scheme is built up from three combinators: transport along a continuous map, the
scheme that repeats one fixed functional, and \texttt{branch}, which asks one
functional and then continues with a scheme chosen arbitrarily from its value. The
last is where adaptivity is created, and the bisection algorithm of
Chapter~\ref{ch:bisection} is built from it.

\chapter{Split covers}
\label{ch:split}

The lower bound rests on a covering notion that measures how far a set is from being
antipodally symmetric.

\begin{definition}[Split half, split cover]
\label{def:split-cover}
\lean{HasSplitCover, IsSplitHalf}
\leanok
A set $P \subseteq E$ is a \emph{split half} if it is closed and disjoint from $-P$;
the associated \emph{split set} is $P \cup (-P)$. A set $K \subseteq E$
\emph{has a split cover with $k$ sets}, written $\mathrm{SC}(K) \le k$, if there are
split halves $P_1, \dots, P_k$ with
\[
  K \subseteq \bigcup_{i=1}^{k} \bigl( P_i \cup (-P_i) \bigr).
\]
\end{definition}

On a split set a sign is available: $+$ on $P$ and $-$ on $-P$. The next lemma is the
elementary engine of the lower bound; applied to the odd part
$\delta(x) = \lambda(x) - \lambda(-x)$ of a measurement, it says that wherever a
measurement separates an antipodal pair, one split set settles the matter.

\begin{lemma}[One nonvanishing odd function]
\label{lem:split-cover-odd}
\lean{hasSplitCover_one_of_odd_ne_zero}
\leanok
\uses{def:split-cover}
Let $W \subseteq E$ be closed and symmetric and let $\delta : E \to \mathbb{R}$ be
continuous and odd with $\delta(x) \neq 0$ for every $x \in W$. Then
$\mathrm{SC}(W) \le 1$.
\end{lemma}
\begin{proof}
\leanok
Take $P = W \cap \{\delta \ge 0\}$. It is closed, and disjoint from $-P$ because
$\delta$ is odd and nonvanishing on $W$.
\end{proof}

\begin{lemma}[Disjoint pieces share their split sets]
\label{lem:split-cover-union}
\lean{HasSplitCover.biUnion_of_pairwise_disjoint}
\leanok
\uses{def:split-cover}
Let $s$ be a finite index set and let $Z_j$, $j \in s$, be compact, symmetric and
pairwise disjoint with $\mathrm{SC}(Z_j) \le h$ for every $j \in s$. Then
$\mathrm{SC}\bigl( \bigcup_{j \in s} Z_j \bigr) \le h$.
\end{lemma}
\begin{proof}
\leanok
\uses{def:split-cover}
Normalise the halves of each $Z_j$ to lie inside $Z_j$ and to be compact, then take
the union over $j$ in each of the $h$ slots. Halves in different slots of different
pieces stay disjoint from each other's reflections because the pieces are disjoint,
compact and symmetric.
\end{proof}

\begin{lemma}[Fattening a split cover]
\label{lem:split-cover-thickening}
\lean{HasSplitCover.exists_cthickening}
\leanok
\uses{def:split-cover}
Let $K \subseteq E$ be compact and symmetric with $\mathrm{SC}(K) \le k$. Then there
is $r > 0$ with $\mathrm{SC}\bigl( \overline{K^r} \bigr) \le k$, where $K^r$ is the
$r$-neighbourhood of $K$.
\end{lemma}
\begin{proof}
\leanok
The halves may be taken compact. A compact set and its reflection are disjoint and
closed, so they stay disjoint after thickening by a small enough radius.
\end{proof}

The next proposition is what converts a covering count into a dimension bound, once
Borsuk--Ulam is available.

\begin{proposition}[A split cover gives an odd map without zeros]
\label{prop:split-cover-odd-map}
\lean{HasSplitCover.exists_odd_continuous_ne_zero}
\leanok
\uses{def:split-cover}
If $\mathrm{SC}(K) \le k$, then there is a continuous odd map
$\Psi : E \to \mathbb{R}^k$ with $\Psi(x) \neq 0$ for every $x \in K$.
\end{proposition}
\begin{proof}
\leanok
Put $\Psi_i(x) = \operatorname{dist}(x, -P_i) - \operatorname{dist}(x, P_i)$. Each
$\Psi_i$ is continuous, and odd because reflection is an isometry. On the split set of
$P_i$ it is nonzero, the two halves being closed and disjoint.
\end{proof}

\chapter{The halving lemma}
\label{ch:halving}

One real value costs a factor of two, and no more. This is the inductive step of the
lower bound.

\begin{lemma}[From a level set to a slab]
\label{lem:split-cover-slab}
\lean{HasSplitCover.exists_pos_inter_preimage_ball}
\leanok
\uses{def:split-cover, lem:split-cover-thickening}
Let $Z \subseteq E$ be compact and symmetric, let $\mu : E \to \mathbb{R}$ be
continuous and even on $Z$, and let $y \in \mathbb{R}$ with
$\mathrm{SC}\bigl( Z \cap \mu^{-1}(y) \bigr) \le h$. Then there is $\rho > 0$ with
$\mathrm{SC}\bigl( Z \cap \mu^{-1}(B(y, \rho)) \bigr) \le h$.
\end{lemma}
\begin{proof}
\leanok
\uses{lem:split-cover-thickening}
Fatten the split cover of the level set to a closed neighbourhood of it. The points of
$Z$ outside that neighbourhood form a compact set whose image under $\mu$ misses $y$,
so their values are bounded away from $y$.
\end{proof}

\begin{theorem}[Halving lemma]
\label{thm:halving}
\lean{hasSplitCover_two_mul_of_forall_fiber}
\leanok
\uses{def:split-cover, lem:split-cover-slab, lem:split-cover-union}
Let $Z \subseteq E$ be compact and symmetric and let $\mu : E \to \mathbb{R}$ be
continuous and even on $Z$. If $\mathrm{SC}\bigl( Z \cap \mu^{-1}(y) \bigr) \le h$ for
every $y \in \mathbb{R}$, then $\mathrm{SC}(Z) \le 2h$.
\end{theorem}
\begin{proof}
\leanok
\uses{lem:split-cover-slab, lem:split-cover-union}
Every value has a ball of values whose slab is still covered by $h$ split sets
(Lemma~\ref{lem:split-cover-slab}). A Lebesgue number for this cover of the compact
set of attained values yields a partition of $\mathbb{R}$ into intervals of that
length; the slabs over the odd-numbered intervals are pairwise disjoint, and so are
those over the even-numbered ones. Lemma~\ref{lem:split-cover-union} covers each of
the two families with $h$ split sets.
\end{proof}

Iterating the halving lemma over the $n$ measurements gives the covering proposition,
which is the whole combinatorial content of the lower bound. Its proof uses no
topology beyond compactness.

\begin{proposition}[Covering proposition]
\label{prop:covering}
\lean{ContMeasurements.AdaptiveScheme.DistinguishesAntipodes.hasSplitCover}
\leanok
\uses{def:distinguishes-antipodes, def:split-cover, thm:halving, lem:split-cover-odd}
Let $K \subseteq E$ be compact and symmetric. If a scheme distinguishes antipodes on
$K$ with $n$ measurements, then $\mathrm{SC}(K) \le 2^n - 1$.
\end{proposition}
\begin{proof}
\leanok
\uses{thm:halving, lem:split-cover-odd}
Induction on $n$. With no measurement, $K$ is empty. For the step, write $\mu$ for the
even part and $\delta$ for the odd part of the first functional. Where $\delta \neq 0$,
one split set suffices (Lemma~\ref{lem:split-cover-odd}). On each level set of $\mu$
where $\delta = 0$, the first measurement gives the same value at $x$ and $-x$, so the
remaining $n - 1$ measurements already distinguish antipodes there, and the inductive
hypothesis covers it with $2^{n-1} - 1$ split sets. The halving lemma turns this into
$2(2^{n-1} - 1)$ for all of $K$ minus the first piece, and
$2(2^{n-1} - 1) + 1 = 2^n - 1$.
\end{proof}

\chapter{The Kuhn triangulation of an integer cube}
\label{ch:grid}

Mathlib contains neither Borsuk--Ulam nor Brouwer's fixed point theorem nor Sperner's
lemma, so the theorem is proved here, combinatorially. The route is Tucker's lemma
\cite{Tucker46} on the triangulation of the cube due to Freudenthal \cite{Freudenthal42}
and Kuhn \cite{Kuhn60}. Describing a simplex by the order relation alone removes every
pseudomanifold argument.

\begin{definition}[Grid and boundary]
\label{def:grid}
\lean{GridSimplex.grid, GridSimplex.IsBoundary}
\leanok
For $n, N \in \mathbb{N}$ the \emph{grid} is
$G = \{ v \in \mathbb{Z}^n : |v_i| \le N \text{ for every } i \}$, and $v \in G$ is a
\emph{boundary point} if $v_c = N$ or $v_c = -N$ for some coordinate $c$.
\end{definition}

\begin{definition}[Labels]
\label{def:label}
\lean{GridSimplex.Label, GridSimplex.negLabel}
\leanok
A \emph{label} is a pair $(i, b)$ with $i$ a coordinate and $b \in \{+, -\}$. The
\emph{opposite} label is $\operatorname{neg}(i, b) = (i, -b)$.
\end{definition}

\begin{definition}[Simplex]
\label{def:simplex}
\lean{GridSimplex.IsSimplex}
\leanok
\uses{def:grid}
A finite set $s \subseteq G$ is a \emph{simplex} if it is nonempty, is a chain for the
coordinatewise order, and satisfies $y_i - x_i \le 1$ for all $x, y \in s$ and every
coordinate $i$.
\end{definition}

These are exactly the faces of the Kuhn triangulation; a top-dimensional one is a
chain $v < v + e_{i_1} < v + e_{i_1} + e_{i_2} < \dots$.

\begin{definition}[Signs and directions]
\label{def:signs}
\lean{GridSimplex.signs, GridSimplex.directions}
\leanok
\uses{def:label}
The label $(i, +)$ is a \emph{sign} of $s$ if $x_i > 0$ for some $x \in s$, and
$(i, -)$ is a sign of $s$ if $x_i < 0$ for some $x \in s$. The \emph{directions} of
$s$ are the coordinates occurring in its signs.
\end{definition}

\begin{lemma}[A simplex lies in one closed orthant]
\label{lem:orthant}
\lean{GridSimplex.IsSimplex.not_mem_signs_negLabel}
\leanok
\uses{def:simplex, def:signs}
For every simplex $s$ and every label $\ell$: if $\ell$ is a sign of $s$, then
$\operatorname{neg} \ell$ is not.
\end{lemma}
\begin{proof}
\leanok
Two vertices with $x_i > 0 > y_i$ would differ by at least $2$ in the coordinate $i$.
\end{proof}

\begin{lemma}[As many signs as directions]
\label{lem:card-signs}
\lean{GridSimplex.IsSimplex.card_signs}
\leanok
\uses{lem:orthant}
For every simplex $s$, the number of signs of $s$ equals the number of directions of
$s$.
\end{lemma}
\begin{proof}
\leanok
\uses{lem:orthant}
By Lemma~\ref{lem:orthant} the coordinate determines the label, so passing from signs
to directions is injective.
\end{proof}

Measured from the least vertex $a$ of a simplex, every vertex $x$ is $a$ raised by one
in the coordinates of its \emph{step set} $\{ i : a_i < x_i \}$; the \emph{rank} of
$x$ is the size of that set. The rank is strictly increasing along the chain, which
bounds the length of a simplex.

\begin{lemma}[Length of a simplex]
\label{lem:card-simplex}
\lean{GridSimplex.IsSimplex.card_le_card_directions_succ}
\leanok
\uses{def:simplex, def:signs}
Every simplex has at most one more vertex than it has directions.
\end{lemma}
\begin{proof}
\leanok
The rank relative to the least vertex is injective on the simplex and takes values
between $0$ and the number of directions.
\end{proof}

Call a simplex \emph{saturated} if it has exactly one more vertex than its total step
set has coordinates. Its ranks then run through every value from $0$ upwards.

\begin{lemma}[Consecutive vertices of a saturated simplex]
\label{lem:unit-step}
\lean{GridSimplex.IsSimplex.exists_eq_add_single}
\leanok
\uses{lem:card-simplex}
In a saturated simplex, two consecutive vertices $x < y$ satisfy $y = x + e_i$ for
exactly one coordinate $i$.
\end{lemma}
\begin{proof}
\leanok
\uses{lem:card-simplex}
Consecutive vertices have consecutive ranks, so the step set of the later one adds
exactly one coordinate.
\end{proof}

The counting argument of Chapter~\ref{ch:tucker} runs over the ways of adding one
vertex to a simplex without adding a sign outside a prescribed set of labels.

\begin{definition}[Extensions]
\label{def:extensions}
\lean{GridSimplex.extensions}
\leanok
\uses{def:simplex, def:signs}
For a simplex $s$ and a set $t$ of labels, the \emph{extensions} of $s$ within $t$ are
the grid points $w \notin s$ such that $s \cup \{w\}$ is a simplex all of whose signs
lie in $t$.
\end{definition}

\begin{definition}[Extremal simplex]
\label{def:extremal}
\lean{GridSimplex.IsExtremal}
\leanok
\uses{def:grid}
A set $s$ of grid points is \emph{extremal} if some coordinate $c$ satisfies $z_c = N$
for every $z \in s$, or $z_c = -N$ for every $z \in s$.
\end{definition}

A simplex can meet the boundary of the cube through a different coordinate at each
vertex without being extremal; extremality is the stronger condition, and it is the
one the counting needs.

\begin{lemma}[No extension]
\label{lem:extensions-empty}
\lean{GridSimplex.extensions_eq_empty}
\leanok
\uses{def:extensions, lem:card-simplex}
A simplex with one more vertex than it has directions has no extension within its own
signs.
\end{lemma}
\begin{proof}
\leanok
\uses{lem:card-simplex}
Such an extension would have two more vertices than directions.
\end{proof}

\begin{lemma}[One extension]
\label{lem:extensions-one}
\lean{GridSimplex.card_extensions_eq_one}
\leanok
\uses{def:extensions, lem:card-simplex, lem:orthant}
Let $N \ge 1$, let $s$ be a simplex with one more vertex than it has directions, and
let $\ell$ be a label such that neither $\ell$ nor $\operatorname{neg} \ell$ is a sign
of $s$. Then $s$ has exactly one extension within $\{\ell\} \cup \operatorname{signs}(s)$.
\end{lemma}
\begin{proof}
\leanok
\uses{lem:card-simplex, lem:unit-step}
The coordinate of $\ell$ vanishes on $s$, and a new vertex must be nonzero there, hence
equal to $\pm 1$ there. It therefore sits above the greatest or below the least vertex
of the chain, and by Lemma~\ref{lem:unit-step} at a single unit step from it. The sign
of $\ell$ decides which end, and $N \ge 1$ keeps the point in the grid.
\end{proof}

\begin{lemma}[Two extensions]
\label{lem:extensions-two}
\lean{GridSimplex.card_extensions_eq_two}
\leanok
\uses{def:extensions, def:extremal, lem:card-simplex}
A simplex with as many vertices as directions and without an extremal coordinate has
exactly two extensions within its own signs.
\end{lemma}
\begin{proof}
\leanok
\uses{lem:unit-step, def:extremal}
Either the total step set already spans all directions, and then the ranks miss a
single value: the two vertices on either side of the gap differ by two unit steps, and
the new vertex takes one of the two intermediate positions. Or one direction is
constant on the simplex, and then the chain grows at either end in that coordinate;
both points stay in the grid precisely because the coordinate is not extremal.
\end{proof}

\begin{lemma}[One extension at the boundary]
\label{lem:extensions-one-extremal}
\lean{GridSimplex.card_extensions_eq_one_of_isExtremal}
\leanok
\uses{def:extensions, def:extremal, lem:card-simplex}
Let $N \ge 1$. A simplex with as many vertices as directions and with an extremal
coordinate has exactly one extension within its own signs.
\end{lemma}
\begin{proof}
\leanok
\uses{lem:extensions-two}
As in Lemma~\ref{lem:extensions-two} the two candidates lie at the two ends of the
chain in the constant direction; the cube blocks the one that would leave the grid.
\end{proof}

\begin{lemma}[Reflection]
\label{lem:simplex-neg}
\lean{GridSimplex.IsSimplex.neg, GridSimplex.signs_neg, GridSimplex.directions_neg}
\leanok
\uses{def:simplex, def:signs}
For every simplex $s$, the reflection $-s$ is a simplex, its signs are the opposites of
the signs of $s$, and its directions are those of $s$.
\end{lemma}
\begin{proof}
\leanok
Reflection reverses the order and negates every coordinate.
\end{proof}

\chapter{Tucker's lemma}
\label{ch:tucker}

Tucker's lemma \cite{Tucker46} is the combinatorial form of the Borsuk--Ulam theorem.
The proof is a parity argument over the simplices on which the labelling is, in the
sense below, complete.

\begin{definition}[Labels on a simplex, happy simplex]
\label{def:happy}
\lean{GridSimplex.labels, GridSimplex.IsHappy}
\leanok
\uses{def:simplex, def:signs}
Let $\lambda$ assign a label to every grid point. The \emph{labels} of a set $s$ are
the values of $\lambda$ on $s$. A simplex $s$ is \emph{happy} if every sign of $s$ is a
label of $s$.
\end{definition}

\begin{definition}[Separating labelling]
\label{def:separates}
\lean{GridSimplex.Separates}
\leanok
\uses{def:simplex, def:label}
A labelling $\lambda$ \emph{separates} if no simplex carries a pair of opposite labels,
that is, $\lambda(x) \neq \operatorname{neg} \lambda(y)$ for all $x, y$ in a common
simplex.
\end{definition}

Tucker's lemma asserts that a labelling antipodal on the boundary never separates.
Everything up to the end of this chapter assumes that one does, and derives a
contradiction.

\begin{lemma}[The three shapes of a happy simplex]
\label{lem:happy-shape}
\lean{GridSimplex.IsHappy.shape}
\leanok
\uses{def:happy, lem:card-signs, lem:card-simplex}
Let $s$ be happy and let $k$ be the number of its directions. Then $s$ has $k+1$
vertices and $k+1$ labels, or $k+1$ vertices and $k$ labels, or $k$ vertices and $k$
labels.
\end{lemma}
\begin{proof}
\leanok
\uses{lem:card-signs, lem:card-simplex}
The signs of $s$ are among its labels, there are as many signs as directions
(Lemma~\ref{lem:card-signs}), a set has at most as many labels as elements, and a
simplex has at most one more vertex than directions (Lemma~\ref{lem:card-simplex}).
\end{proof}

\begin{definition}[Door, degree]
\label{def:door}
\lean{GridSimplex.IsDoor, GridSimplex.deg}
\leanok
\uses{def:happy}
A happy simplex $\sigma$ is a \emph{door} of a happy simplex $\tau$ if $\sigma$ is
$\tau$ with one vertex deleted and every sign of $\tau$ is a label of $\sigma$. The
\emph{degree} of a happy simplex is the number of doors it has, counted from both
sides: the happy simplices of which it is a door, plus its own doors.
\end{definition}

\begin{lemma}[Counting the doors twice]
\label{lem:door-count}
\lean{GridSimplex.sum_upDeg_eq_sum_downDeg}
\leanok
\uses{def:door}
Summed over all happy simplices, the number of doors above equals the number of doors
below.
\end{lemma}
\begin{proof}
\leanok
Both sums count the pairs $(\sigma, \tau)$ with $\sigma$ a door of $\tau$, in the two
possible orders.
\end{proof}

\begin{lemma}[An even number of odd degrees]
\label{lem:even-odd-degree}
\lean{GridSimplex.even_card_odd_deg}
\leanok
\uses{def:door, lem:door-count}
The number of happy simplices of odd degree is even.
\end{lemma}
\begin{proof}
\leanok
\uses{lem:door-count}
By Lemma~\ref{lem:door-count} the total degree is twice the number of doors, hence
even, and a sum of natural numbers is even exactly when an even number of its terms
are odd.
\end{proof}

\begin{lemma}[A fixed-point-free involution]
\label{lem:involution}
\lean{Finset.even_card_of_involutive}
\leanok
A finite set carrying an involution without fixed points has an even number of
elements.
\end{lemma}
\begin{proof}
\leanok
Induction on the cardinality: remove a point and its image, which are distinct.
\end{proof}

\begin{theorem}[Which happy simplices have odd degree]
\label{thm:odd-degree}
\lean{GridSimplex.odd_deg_iff}
\leanok
\uses{def:door, def:separates, def:extremal, lem:happy-shape, lem:extensions-empty,
lem:extensions-one, lem:extensions-two, lem:extensions-one-extremal}
Let $N \ge 1$ and let $\lambda$ separate. A happy simplex has odd degree if and only if
it is the origin alone, or it has as many vertices as directions and an extremal
coordinate.
\end{theorem}
\begin{proof}
\leanok
\uses{lem:happy-shape, lem:extensions-empty, lem:extensions-one, lem:extensions-two,
lem:extensions-one-extremal}
Take the three shapes of Lemma~\ref{lem:happy-shape} in turn. With $k+1$ vertices and
$k+1$ labels, the labels are distinct and exceed the signs by exactly one label
$\ell$; deleting the vertex carrying $\ell$ is the only door below, and since $\lambda$
separates, $\operatorname{neg}\ell$ is not a sign either, so
Lemma~\ref{lem:extensions-one} gives exactly one door above. The degree is $2$, except
for the simplex $\{0\}$, which has no door below. With $k+1$ vertices and $k$ labels,
the labels are the signs and exactly one label is repeated: the two vertices carrying
it give two doors below, and Lemma~\ref{lem:extensions-empty} leaves none above. With
$k$ vertices and $k$ labels, the labels are again the signs, so no vertex may be
deleted, and Lemmas~\ref{lem:extensions-two} and~\ref{lem:extensions-one-extremal}
give two doors above, or one when a coordinate is extremal.
\end{proof}

\begin{theorem}[Tucker's lemma]
\label{thm:tucker}
\lean{GridSimplex.tucker}
\leanok
\uses{def:grid, def:label, def:separates, thm:odd-degree, lem:even-odd-degree,
lem:involution, lem:simplex-neg}
Let $n \ge 0$ and $N \ge 1$, and let $\lambda$ assign a label to every point of the
grid $G \subseteq \mathbb{Z}^n$ so that
$\lambda(-v) = \operatorname{neg} \lambda(v)$ for every boundary point $v$. Then there
are $u, w \in G$ with $u \le w$, with $w_i - u_i \le 1$ for every coordinate $i$, and
with $\lambda(u) = \operatorname{neg} \lambda(w)$.
\end{theorem}
\begin{proof}
\leanok
\uses{thm:odd-degree, lem:even-odd-degree, lem:involution, lem:simplex-neg}
Suppose not. Any two points of a simplex are comparable and differ by at most one in
each coordinate, so $\lambda$ separates, and Theorem~\ref{thm:odd-degree} applies. The
happy simplices of odd degree are then the origin alone together with the extremal
ones. An extremal simplex has all its vertices on the boundary, so the labelling is
antipodal there, and reflection maps the extremal happy simplices to themselves
(Lemma~\ref{lem:simplex-neg}); it has no fixed point, since the extremal coordinate
would have to take both values $N$ and $-N$. By Lemma~\ref{lem:involution} their number
is even, so the number of happy simplices of odd degree is odd, contradicting
Lemma~\ref{lem:even-odd-degree}.
\end{proof}

\chapter{The Borsuk--Ulam theorem}
\label{ch:borsuk}

The odd-map form of the Borsuk--Ulam theorem \cite{Borsuk33} states that a continuous
odd map from $S^{m-1}$ into $\mathbb{R}^k$ with $k < m$ vanishes somewhere. It is
reached from Tucker's lemma in three steps: the cube, the sphere of the sup norm, and
the Euclidean sphere.

In this chapter $\mathbb{R}^n$ carries the sup norm, so the closed unit ball is the
cube $[-1, 1]^n$ and the unit sphere is its boundary.

\begin{definition}[Label of a point]
\label{def:max-label}
\lean{CubeLabel.label, CubeLabel.maxIdx}
\leanok
\uses{def:label}
For $y \in \mathbb{R}^n$ with $n \ge 1$, let $i(y)$ be the least coordinate with
$|y_{i(y)}| = \|y\|$. The \emph{label} of $y$ is $\bigl( i(y), \operatorname{sign}
y_{i(y)} \bigr)$.
\end{definition}

Taking the least such coordinate is what makes the labelling antipodal: reflection
reverses every coordinate without moving the set of largest ones.

\begin{lemma}[Reflection reverses the label]
\label{lem:label-neg}
\lean{CubeLabel.label_neg}
\leanok
\uses{def:max-label}
For every $y \neq 0$, the label of $-y$ is the opposite of the label of $y$.
\end{lemma}
\begin{proof}
\leanok
The coordinates attaining the sup norm are the same for $y$ and $-y$, so the chosen
coordinate is the same, and the sign there is reversed.
\end{proof}

\begin{theorem}[Borsuk's antipodal theorem for the cube]
\label{thm:borsuk-cube}
\lean{borsukUlam_cube}
\leanok
\uses{thm:tucker, def:max-label, lem:label-neg}
Let $n \ge 1$ and let $g : \mathbb{R}^n \to \mathbb{R}^n$ be continuous on the cube
$[-1, 1]^n$ and odd on its boundary. Then $g(x) = 0$ for some $x$ in the cube.
\end{theorem}
\begin{proof}
\leanok
\uses{thm:tucker, lem:label-neg}
Suppose $g$ never vanishes on the cube. Then $\varepsilon = \min_{[-1,1]^n} \|g\| > 0$.
Choose a mesh $1/N$ below the modulus of continuity of $g$ for $2\varepsilon$ and label
each point $v$ of the grid $[-N, N]^n$ by the label of $g(v/N)$. On the boundary of the
grid $\|v/N\| = 1$, so $g$ is odd there and the labelling is antipodal
(Lemma~\ref{lem:label-neg}). Tucker's lemma gives two grid points $u, w$ with opposite
labels and $\|u/N - w/N\| \le 1/N$. Their images then differ by at least $2\varepsilon$
in one coordinate, although the points are closer than the mesh.
\end{proof}

The unit sphere of the sup norm in $\mathbb{R}^{n+1}$ is reached from the cube
$[-1, 1]^n$ by adding the slack of the sup norm as a further coordinate and projecting
radially.

\begin{definition}[Lift and radial projection]
\label{def:lift}
\lean{CubeSphere.lift, CubeSphere.proj}
\leanok
For $x \in \mathbb{R}^n$ put $\operatorname{lift}(x) = (x, 1 - \|x\|) \in
\mathbb{R}^{n+1}$ and
$\operatorname{proj}(x) = \operatorname{lift}(x) / \|\operatorname{lift}(x)\|$.
\end{definition}

\begin{lemma}[The lift stays away from the origin]
\label{lem:lift-pos}
\lean{CubeSphere.half_le_norm_lift}
\leanok
\uses{def:lift}
For every $x \in \mathbb{R}^n$, $\|\operatorname{lift}(x)\| \ge 1/2$.
\end{lemma}
\begin{proof}
\leanok
Either $\|x\| \ge 1/2$, and the lift is at least as long, or $\|x\| < 1/2$ and the last
coordinate exceeds $1/2$.
\end{proof}

So $\operatorname{proj}$ is continuous on all of $\mathbb{R}^n$ and takes values in the
unit sphere; on the boundary of the cube the slack vanishes, so it is odd there.

\begin{theorem}[Borsuk--Ulam for the sup norm]
\label{thm:borsuk-pi}
\lean{borsukUlam_pi}
\leanok
\uses{thm:borsuk-cube, def:lift, lem:lift-pos}
Let $k < m$ and let $f : \mathbb{R}^m \to \mathbb{R}^k$ be odd and continuous on the
unit sphere of the sup norm. Then $f(x) = 0$ for some $x$ of that sphere.
\end{theorem}
\begin{proof}
\leanok
\uses{thm:borsuk-cube, lem:lift-pos}
For $k = 0$ any point of the sphere will do. Otherwise write $m = n + 1$ with
$n \ge k$, pad the values of $f$ with zeros to $n$ coordinates, and apply
Theorem~\ref{thm:borsuk-cube} to the composition with $\operatorname{proj}$, which is
continuous on the cube and odd on its boundary.
\end{proof}

\begin{theorem}[Borsuk--Ulam theorem]
\label{thm:borsuk-ulam}
\lean{borsukUlam}
\leanok
\uses{thm:borsuk-pi}
Let $k < m$ and let $f : \mathbb{R}^m \to \mathbb{R}^k$ be odd and continuous on the
Euclidean unit sphere. Then $f(x) = 0$ for some $x$ of that sphere.
\end{theorem}
\begin{proof}
\leanok
\uses{thm:borsuk-pi}
The map $u \mapsto u / \|u\|_2$ sends the sphere of the sup norm to the Euclidean one,
continuously and oddly. Compose and apply Theorem~\ref{thm:borsuk-pi}.
\end{proof}

The form the lower bound consumes drops both the Euclidean norm and the radius. No
property of the norm is used beyond homogeneity, which is why the lower bound holds for
any norm.

\begin{corollary}[Borsuk--Ulam for an arbitrary norm and radius]
\label{cor:borsuk-ulam-norm}
\lean{Function.Odd.exists_norm_eq_and_eq_zero_of_lt_finrank}
\leanok
\uses{thm:borsuk-ulam}
Let $E$ be a real normed space with $\dim E = m$, let $k < m$, let $R > 0$ and let
$\Psi : E \to \mathbb{R}^k$ be continuous and odd. Then there is $x \in E$ with
$\|x\| = R$ and $\Psi(x) = 0$.
\end{corollary}
\begin{proof}
\leanok
\uses{thm:borsuk-ulam}
A linear homeomorphism carries the Euclidean unit sphere onto a compact symmetric set,
and $u \mapsto (R / \|e(u)\|) \, e(u)$ sends that set onto the sphere of radius $R$,
oddly and continuously.
\end{proof}

\chapter{The lower bound}
\label{ch:lower}

The lower bound is due to Ullrich \cite{U26}. Combining the covering proposition with
Borsuk--Ulam turns a count of split sets into a bound on the dimension.

\begin{theorem}[Antipodal coincidence]
\label{thm:antipodal-coincidence}
\lean{ContMeasurements.exists_norm_eq_and_info_eq_info_neg}
\leanok
\uses{def:info, prop:covering, prop:split-cover-odd-map, cor:borsuk-ulam-norm}
Let $E$ be a real normed space with $\dim E = m$, let $S$ be an adaptive scheme on $E$,
let $n \in \mathbb{N}_0$ with $2^n \le m$, and let $R > 0$. Then there is $x \in E$ with
$\|x\| = R$ and $N_n(x) = N_n(-x)$.
\end{theorem}
\begin{proof}
\leanok
\uses{prop:covering, prop:split-cover-odd-map, cor:borsuk-ulam-norm}
If not, the scheme distinguishes antipodes on the sphere of radius $R$, which is
compact and symmetric, so Proposition~\ref{prop:covering} covers that sphere with
$2^n - 1$ split sets. Proposition~\ref{prop:split-cover-odd-map} turns the cover into a
continuous odd map into $\mathbb{R}^{2^n - 1}$ without zeros on the sphere. Since
$2^n - 1 < m$, Corollary~\ref{cor:borsuk-ulam-norm} produces a zero of norm $R$.
\end{proof}

\begin{corollary}[Counting form]
\label{cor:counting}
\lean{ContMeasurements.AdaptiveScheme.DistinguishesAntipodes.finrank_add_one_le_two_pow}
\leanok
\uses{def:distinguishes-antipodes, thm:antipodal-coincidence}
Let $E$ be a finite-dimensional real normed space, $R > 0$, and let a scheme
distinguish antipodes on the sphere of radius $R$ with $n$ measurements. Then
$\dim E + 1 \le 2^n$.
\end{corollary}
\begin{proof}
\leanok
\uses{thm:antipodal-coincidence}
Otherwise $2^n \le \dim E$, and Theorem~\ref{thm:antipodal-coincidence} produces an
antipodal pair on the sphere with the same information.
\end{proof}

\begin{corollary}[Below the threshold nothing is achieved]
\label{cor:lower-error}
\lean{ContMeasurements.exists_norm_eq_and_le_norm_sub}
\leanok
\uses{thm:antipodal-coincidence}
Let $E$ be a real normed space with $\dim E = m$, let $S$ be an adaptive scheme on $E$,
let $n \in \mathbb{N}_0$ with $2^n \le m$, let $R > 0$ and let
$\Phi : \mathbb{R}^n \to E$ be arbitrary. Then there is $x \in E$ with $\|x\| = R$ and
$R \le \|x - \Phi(N_n(x))\|$.
\end{corollary}
\begin{proof}
\leanok
\uses{thm:antipodal-coincidence}
Take an antipodal pair with the same information. Both points receive the same output,
and they are $2R$ apart, so one of them is at distance at least $R$ from it.
\end{proof}

\chapter{Coloured covers}
\label{ch:colored}

The upper bound is due to Krieg, Novak and Ullrich \cite{KNU26}. Its combinatorial
input is a covering of the space by small pieces whose colours can be told apart
cheaply.

\begin{definition}[Coloured cover]
\label{def:colored-cover}
\lean{ColoredCover}
\leanok
A \emph{coloured cover} of a metric space $X$ with $K$ colours, diameter $\varepsilon$
and separation $c$ consists of pieces $A_i$ covering $X$ and a colour
$\kappa(i) \in \{1, \dots, K\}$ for each piece, such that
$\operatorname{dist}(x, y) \le \varepsilon$ for all $x, y$ in one piece, and
$\operatorname{dist}(x, y) \ge c$ whenever $x \in A_i$, $y \in A_j$ with $i \neq j$ and
$\kappa(i) = \kappa(j)$.
\end{definition}

The \emph{colour class} of a colour is the union of the pieces carrying it. Within one
colour class the pieces are $c$-separated, which is what makes a single measurement
enough to name the piece.

\begin{lemma}[Rescaling]
\label{lem:scale}
\lean{ColoredCover.scale}
\leanok
\uses{def:colored-cover}
For every $s > 0$, a coloured cover with $K$ colours, diameter $\varepsilon$ and
separation $c$ yields one with $K$ colours, diameter $s\varepsilon$ and separation
$sc$.
\end{lemma}
\begin{proof}
\leanok
Replace each piece $A$ by $sA$.
\end{proof}

\begin{theorem}[The Lebesgue cover of $\mathbb{R}^m$]
\label{thm:lebesgue-cover}
\lean{ColoredCover.lebesgue}
\leanok
\uses{def:colored-cover}
The space $\mathbb{R}^m$ with the maximum norm has a coloured cover with $m + 1$
colours, diameter $3/2$ and a positive separation.
\end{theorem}
\begin{proof}
\leanok
The pieces are indexed by an integer point together with a subset of the coordinates,
and are neighbourhoods of the faces of the integer grid; the colour is the size of the
subset. Two distinct pieces of the same colour meet different faces of the same
dimension and are therefore separated.
\end{proof}

The number $m + 1$ is the Lebesgue covering dimension bound and cannot be lowered.

\begin{definition}[Logistic function and logit]
\label{def:logistic}
\lean{Real.logistic, Real.logit}
\leanok
The \emph{logistic function} is $\ell(s) = e^s / (1 + e^s)$ and the \emph{logit} is
$\ell^{-1}(p) = \log\bigl( p / (1 - p) \bigr)$.
\end{definition}

The logistic function is a continuous injection $\mathbb{R} \to (0,1)$; it is what
compresses a countable family of values into a bounded range in the identification
measurement below.

\chapter{Bisection and identification}
\label{ch:bisection}

Two ingredients turn a coloured cover into an algorithm: one measurement names the
piece within a colour class, and $\lceil \log_2 K \rceil$ adaptive measurements find
the colour.

\begin{theorem}[The identification measurement]
\label{thm:identification}
\lean{ContMeasurements.exists_identification}
\leanok
\uses{def:colored-cover, def:logistic}
Let $X$ be a nonempty metric space, let $cc$ be a coloured cover of $X$ with countable
index set, diameter $\varepsilon$ and separation $c > 0$, and let $r$ be a colour.
Then there are a continuous functional $\varphi : X \to \mathbb{R}$ and a map
$\mathrm{dec} : \mathbb{Z} \to X$ such that for every $x$ in the closure of the colour
class of $r$ there is $n \in \mathbb{Z}$ with $\varphi(x) = n$ and
$\operatorname{dist}\bigl( x, \mathrm{dec}(n) \bigr) \le \varepsilon + c/8$.
\end{theorem}
\begin{proof}
\leanok
\uses{def:logistic}
Enumerate the pieces of the colour class and attach to the $j$-th piece a collar
function, supported in its $c/8$-neighbourhood and equal to $1$ on the piece. The
functional is the sum of the collar functions weighted by the logit of an injective
encoding of $j$; the weights make the sum locally finite and continuous, and its value
on a point of the colour class names the piece.
\end{proof}

\begin{theorem}[Bisection over colours]
\label{thm:bisection}
\lean{ContMeasurements.exists_bisection_scheme}
\leanok
\uses{def:adaptive-scheme, def:info}
Let $X$ be a nonempty metric space, let $A_1, \dots, A_{K}$ be sets, and suppose that
for each $r$ there are a scheme of $\ell$ measurements and a reconstruction map that is
$\eta$-accurate on $\overline{A_r}$. Then for every $L \in \mathbb{N}_0$ and every set
$C$ of at most $2^L$ colours there are a scheme of $\ell + L$ measurements and a
reconstruction map that is $\eta$-accurate on $\overline{\bigcup_{r \in C} A_r}$.
\end{theorem}
\begin{proof}
\leanok
\uses{def:adaptive-scheme}
Induction on $L$. Split $C$ into two halves of at most $2^{L-1}$ colours, take the
schemes the inductive hypothesis provides for each half, and put them together with the
branch combinator, asking first a functional that separates the two halves.
\end{proof}

\chapter{The upper bound and the two bounds together}
\label{ch:upper}

\begin{theorem}[Upper bound]
\label{thm:upper-bound}
\lean{ContMeasurements.exists_recovery_scheme}
\leanok
\uses{def:adaptive-scheme, def:info, thm:lebesgue-cover, lem:scale, thm:identification,
thm:bisection}
Let $E$ be a real normed space with $\dim E = d \ge 1$ and let $\varepsilon > 0$. Then
there are an adaptive scheme $S$ on $E$ and a map
$\Phi : \mathbb{R}^{\lceil \log_2 d \rceil + 1} \to E$ with
\[
  \| x - \Phi( N_{\lceil \log_2 d \rceil + 1}(x) ) \| \le \varepsilon
  \qquad \text{for every } x \in E .
\]
\end{theorem}
\begin{proof}
\leanok
\uses{thm:lebesgue-cover, lem:scale, thm:identification, thm:bisection}
In coordinates: rescale the Lebesgue cover of $\mathbb{R}^d$
(Theorem~\ref{thm:lebesgue-cover}, Lemma~\ref{lem:scale}) to diameter below
$\varepsilon$. One identification measurement is $\varepsilon$-accurate on each of the
$d + 1$ colour classes (Theorem~\ref{thm:identification}), and
$\lceil \log_2 (d+1) \rceil$ bisection measurements find the colour
(Theorem~\ref{thm:bisection}), which is
$\lceil \log_2 d \rceil + 1$ measurements in total. For a general norm, transport the
coordinate algorithm along a linear homeomorphism and rescale the accuracy by the norm
of its inverse.
\end{proof}

\begin{definition}[Measurement complexity]
\label{def:min-measurements}
\lean{ContMeasurements.minMeasurements}
\leanok
\uses{def:info}
For a normed space $E$ and $R, \varepsilon > 0$, let $n^*(E, R, \varepsilon)$ be the
least $n$ for which there are an adaptive scheme $S$ on $E$ and a map
$\Phi : \mathbb{R}^n \to E$ with
$\|x - \Phi(N_n(x))\| \le \varepsilon$ for every $x$ of the closed ball of radius $R$.
\end{definition}

\begin{corollary}[The two bounds]
\label{cor:sandwich}
\lean{ContMeasurements.clog_le_minMeasurements_and_le,
ContMeasurements.finrank_add_one_le_two_pow_minMeasurements}
\leanok
\uses{def:min-measurements, thm:upper-bound, cor:lower-error}
Let $E$ be a finite-dimensional real normed space with $\dim E = d \ge 1$ and let
$0 < \varepsilon < R$. Then $d + 1 \le 2^{\,n^*(E, R, \varepsilon)}$ and
\[
  \lceil \log_2 (d+1) \rceil \;\le\; n^*(E, R, \varepsilon)
  \;\le\; \lceil \log_2 d \rceil + 1 .
\]
\end{corollary}
\begin{proof}
\leanok
\uses{thm:upper-bound, cor:lower-error}
The upper bound is Theorem~\ref{thm:upper-bound}. For the lower bound, an algorithm
accurate to within $\varepsilon < R$ on the ball of radius $R$ cannot err by $R$ on the
sphere of radius $R$, so Corollary~\ref{cor:lower-error} forbids
$2^{n} \le d$.
\end{proof}

\begin{corollary}[Equality in dimension a power of two]
\label{cor:power-of-two}
\lean{ContMeasurements.minMeasurements_eq_of_finrank_eq_two_pow}
\leanok
\uses{cor:sandwich}
Let $E$ be a real normed space with $\dim E = 2^k$ and let $0 < \varepsilon < R$. Then
$n^*(E, R, \varepsilon) = k + 1$.
\end{corollary}
\begin{proof}
\leanok
\uses{cor:sandwich}
Both bounds of Corollary~\ref{cor:sandwich} equal $k + 1$.
\end{proof}

The two bounds meet when $d$ is a power of two and differ by one otherwise.

\bpbibheading
\begin{thebibliography}{99}

\bibitem{Borsuk33}
K.~Borsuk,
\emph{Drei S\"atze \"uber die $n$-dimensionale euklidische Sph\"are},
Fund.\ Math.\ \textbf{20} (1933), 177--190.

\bibitem{Freudenthal42}
H.~Freudenthal,
\emph{Simplizialzerlegungen von beschr\"ankter Flachheit},
Ann.\ of Math.\ (2) \textbf{43} (1942), 580--582.

\bibitem{KNU26}
D.~Krieg, E.~Novak, M.~Ullrich,
\emph{How many questions do you need to find a vector?},
preprint, 2026.

\bibitem{Kuhn60}
H.~W.~Kuhn,
\emph{Some combinatorial lemmas in topology},
IBM J.\ Res.\ Develop.\ \textbf{4} (1960), 518--524.

\bibitem{Tucker46}
A.~W.~Tucker,
\emph{Some topological properties of disk and sphere},
in: Proc.\ First Canadian Math.\ Congress (Montreal, 1945),
University of Toronto Press, 1946, 285--309.

\bibitem{U26}
M.~Ullrich,
\emph{Optimal learning based on adaptive continuous measurements},
preprint, 2026.

\end{thebibliography}


\end{document}
